\section{Power Cable Sound Emission --- Summary of the On-Site Measurements}
\label{sec:onsite_summary}

This section summarizes the main results obtained from the \textit{on-field}
analysis.

\begin{itemize}
	\item[--] \textbf{Description of the problem}:
	\begin{itemize}
		\item[--] The Les Arcs ropeway plant comprises two drive sections, each
		consisting of a PSM motor and four electrical drives connected,
		respectively, to eight galvanically isolated three-phase systems.
		\item[--] Each machine system is connected to its electrical drive by a
		shielded three-phase power cable with a symmetrical inner ground
		conductor, as shown in Figure~\ref{electrical_power_cable}.
		\item[--] The motor power cable of section~1 is \SI{40}{\meter} long.
		\item[--] The motor power cable of section~2 is \SI{120}{\meter} long.
		\item[--] The \SI{120}{\meter} section exhibits audible sound emission from
		the cables themselves; the \SI{40}{\meter} section does not.
	\end{itemize}

	\item[--] \textbf{On-site measurement results}:
	\begin{itemize}
		\item[--] The magnitude of the audible sound emission is proportional to
		the plant load.
		\item[--] The tone of the audible sound emission is correlated to the
		inverter PWM switching frequency.
		\item[--] With a constant PWM carrier frequency at the inverter side, the
		audible emission is narrowband, centered on a single tone.
		\item[--] With random PWM enabled at the inverter side, the emission becomes
		broadband (spread-spectrum), and the resulting sound is subjectively more
		acceptable.
		\item[--] The common-mode (zero-sequence) current is
		$I_z^{\mathrm{inv}}=\SI{8}{\ampere}$ rms at the drive output and
		$I_z^{\mathrm{m}}=\SI{2.5}{\ampere}$ rms at the machine input. The larger
		part of the common-mode current therefore returns through the cable shield
		rather than through the machine, which bounds the fraction available to
		drive bearing current. This common-mode path concerns bearing-current and
		EMC behavior and is distinct from the acoustic mechanism discussed below,
		which is a differential-mode effect.
	\end{itemize}

	\item[--] \textbf{Preliminary conclusions}:
	\begin{itemize}
		\item[--] The power cable on both sections experiences internal inter-phase
		forces proportional to the product of the fundamental phase current $I_1$
		and the switching-frequency current ripple of the adjacent phase
		$\hat{\tilde\imath}_{sw}$:
		\begin{equation}
			\boxed{\;\hat F_{sw}\approx\frac{\mu_0}{2\pi d}\,\hat I_1\,\hat{\tilde\imath}_{sw}\;}
			\label{differential_force}
		\end{equation}
		acting at a frequency $\approx f_{sw}$ with $\pm f_1$ sidebands.
		\item[--] The force generated by \eqref{differential_force} is distributed
		along the whole cable length.
		\item[--] The power cable can be modelled as a mixed string--beam system,
		described by the partial differential equation
		\begin{equation}
			\rho_A\frac{\partial^2 w}{\partial t^2}
			+EI\frac{\partial^4 w}{\partial z^4}
			-T\frac{\partial^2 w}{\partial z^2}
			=f(z,t).
			\label{string_beam_model}
		\end{equation}
		\item[--] The per-unit-length electromagnetic force is approximately equal
		on both sections, under the assumption of the same fundamental current
		$I_1$, conductor spacing $d$, and switching ripple
		$\hat{\tilde\imath}_{sw}$. The difference in emission therefore does not
		originate in the excitation, but in the mechanical--acoustic response of
		the two runs.
		\item[--] The string--beam system~\eqref{string_beam_model} has a dense
		spectrum of span resonances whose modal density scales with cable length.
		On the \SI{120}{\meter} run a force harmonic at $f_{sw}\pm m f_1$ is far
		more likely to coincide with a lightly damped span resonance and be
		resonantly amplified; on the \SI{40}{\meter} run the sparser resonances can
		fall between the force lines, leaving the response near the non-resonant
		level. This resonant coincidence is the dominant reason why the
		\SI{120}{\meter} section emits audible sound while the \SI{40}{\meter}
		section does not.
		\item[--] The larger radiating length of the \SI{120}{\meter} cable adds a
		secondary contribution of about $10\log_{10}(120/40)\approx\SI{4.8}{\decibel}$
		for equal surface velocity, but this alone cannot account for the
		audible/inaudible contrast between the two sections.
		\item[--] The resonant-coupling hypothesis can be confirmed on site by an
		impact (bump) test on a \SI{120}{\meter} span to identify its mechanical
		natural frequencies, and by verifying their coincidence with the force
		harmonics $f_{sw}\pm m f_1$.
	\end{itemize}
\end{itemize}
